\chapter{研究架構}
\label{ch:framework}

\section{概念架構}

\method{} 將控制問題拆成慢速全域規劃與快速局部修正。慢速端由 frozen \smolvla{} 根據任務語言、影像與狀態產生 base action chunk；快速端在每一個 control tick 讀取當下 wrist patch features、robot state、chunk position 與 selected base action，輸出 7 維 residual correction。部署時只執行有界 residual merge。

\begin{figure}[htbp]
  \centering
  \includegraphics[width=0.98\textwidth]{fig1_solution_schematic.png}
  \caption{\method{} 架構圖。Frozen \smolvla{} slow planner 依 SmolVLA input contract 接收 language、top-view RGB、wrist RGB 與 robot state，經 VLM/action-expert stack 以低頻產生 base action chunks 與 cached slow context；frozen DINOv3 wrist encoder 與 trainable fast residual module 使用 current wrist feedback、state、chunk context 與 selected base action 預測有界 residual。Final action 由 residual 經 clip 與 scale 後加到 selected base action 形成。}
  \label{fig:solution_schematic}
\end{figure}

此設計的關鍵不是讓 fast module 重新規劃任務，而是讓它在 slow planner 的 nominal action 附近做 bounded correction。因此本文不能把 \method{} 描述為 wrist-only controller；準確描述是：\textbf{高頻新視覺證據來自腕部相機，而全域任務語意仍由凍結 slow planner 提供。}

\section{問題變數}

令 slow planner 在時間 $t$ 產生 action chunk：

\begin{equation}
  A_t=(a_t^0,a_t^1,\ldots,a_t^{H-1}).
\end{equation}

在第 $k$ 個 chunk position 執行時，base action 為 $\abase=a_t^k$。若 robot 與環境已經偏離 $A_t$ 生成時的 observation，則 $\abase$ 可能不再是最佳 action。\method{} 以 fast residual module $f_\theta$ 預測：

\begin{equation}
  \deltah=f_\theta(x_t^{wrist},s_t,\abase,k,c_t^{slow}),
\end{equation}

其中 $x_t^{wrist}$ 是 wrist DINO patch tokens，$s_t$ 是 robot state，$c_t^{slow}$ 是 cached slow-planner context。最後執行：

\begin{equation}
  \afinal=\abase+\alpha\cdot\clip(\deltah,-\dmax,\dmax).
  \label{eq:merge}
\end{equation}

\section{Claim boundary}

本文把 evidence 分成四層：

\begin{enumerate}
  \item \textbf{主結果：} LIBERO-Spatial synchronous plan=50 execution/replan sweep 與 matched chunk sweep。
  \item \textbf{診斷結果：} residual alpha/clip calibration 與 async-timestep planner-delay stress test。
  \item \textbf{研究流程紀錄：} dataset recording、cache construction、artifact packaging 與 registry 管理。
  \item \textbf{未完成驗證：} no-residual-clip controls、multi-seed、causal visual ablations、A2C2 same-backbone comparison 與 real-robot validation。
\end{enumerate}

這樣分層可避免把單一 seed simulation 結果寫成過強 claim，也符合本論文作為碩士論文「研究流程完整紀錄」的定位。
